% ======================================================================
% SEÇÃO 4 — RESULTADOS
% ======================================================================
\section{Resultados}
\label{sec:res}

Os resultados reportados correspondem a execu\c{c}\~oes reais das su\'ites de
testes automatizados realizadas no ambiente de desenvolvimento do projeto. A
Tabela~\ref{tab:testes} consolida os resultados por camada.

\begin{table}[htbp]
\centering
\caption{Resultado da execu\c{c}\~ao das su\'ites de testes automatizados.}
\label{tab:testes}
\begin{tabular}{lcccc}
\toprule
\textbf{Camada} & \textbf{Ferramenta} & \textbf{Arquivos} & \textbf{Testes} & \textbf{Resultado} \\
\midrule
Backend (Python) & pytest & 8 & 40 & 40 aprovados \\
Firmware (\daq) & PlatformIO (host) & 1 & 14 & 14 aprovados \\
IHM Web & Vitest & 4 & 13 & 13 aprovados \\
\midrule
\textbf{Total} & --- & 13 & 67 & 67 aprovados \\
\bottomrule
\end{tabular}
\end{table}

\subsection{Testes de l\'ogica e integra\c{c}\~ao}

Os testes do backend validaram as transi\c{c}\~oes da FSM e a matriz de
acionamento em cada fase; o c\'alculo da sa\'ida do PID, a satura\c{c}\~ao e o
anti-windup; o c\'alculo da taxa de aquecimento e a sele\c{c}\~ao das duas
regi\~oes da rampa (raz\~ao de taxas abaixo de $0\,\celsius$ e PID acima); a
escrita/leitura at\^omicas da persist\^encia com restaura\c{c}\~ao do backup; o
codec JSON do enlace serial; e os endpoints da API. Os testes do firmware
validaram o parser de pacotes JSON, a rejei\c{c}\~ao de JSON inv\'alido, a
constru\c{c}\~ao do pacote de resposta, o watchdog (disparo, rearme e
configura\c{c}\~ao do tempo limite) e o pulso de toggle da bomba. Os testes do
frontend validaram o consumo da telemetria WebSocket, o sincronismo dos controles
manuais com a telemetria e a persist\^encia da prefer\^encia de tema.

\subsection{Testes funcionais e de seguran\c{c}a}

Os testes funcionais exercitaram o ciclo autom\'atico $T_0 \to T_3$ com o
simulador do \daq, confirmando que: (i) o ciclo somente inicia a partir do estado
seguro; (ii) as dura\c{c}\~oes das fases seguem os par\^ametros configurados; e
(iii) a telemetria publicada cont\'em o estado, o progresso da etapa e o
progresso geral do ciclo, habilitando o Diagrama de Tempos da IHM. Os cen\'arios
de seguran\c{c}a simularam a perda de comunica\c{c}\~ao serial e confirmaram as
duas camadas de prote\c{c}\~ao: o watchdog do firmware desliga os fornos na
aus\^encia de pacotes v\'alidos por mais de 1~s, e o watchdog do backend comuta a
FSM para o estado seguro quando o \daq deixa de responder por mais de 1~s,
atendendo ao crit\'erio de fornos desligados em $\le 1$~s.

A Figura~\ref{fig:ihm} apresenta a IHM Web em opera\c{c}\~ao, exibindo o Diagrama
de Tempos, o progresso da etapa, os controles e atuadores e os gr\'aficos de
tend\^encia por forno.

\begin{figure}[htbp]
\centering
\includegraphics[width=0.9\linewidth]{frontendweb.png}
\caption{IHM Web em modo de monitoramento.}
\label{fig:ihm}
\par\vspace{2pt}{\footnotesize Fonte: captura de tela da IHM do projeto.}
\end{figure}
